\section*{Bribery, Corruption, and Corporations}
The apparently recent emphasis of ethics in engineering arises from recognition and collective desire to combat bribery and corruption in engineering projects.   The consensus that bribery and corruption is unacceptable is not just based on moral arguments but is thought to consume about 3\% of the gross domestic product of the entire world, about 5 trillion dollars globally or about \$900 per person on the planet\footnote{ http://en.wikipedia.org/wiki/Image:World GDP Capita1-2003 A.D.png ; I assumed 6 billion world population}. 

Figure~\ref{fig:soul_for_a_donut} is a humorous view of the moral transaction of corruption and bribery.  Wally being of little character sells his soul for a donut!

\begin{figure}[h!] %  figure placement: here, top, bottom, or page
   \centering
   \includegraphics[width=5in]{./figures/soul_for_a_donut.pdf} 
   \caption{Wally sells his soul for a donut; corruption costs a lot more than donuts, but is essentially the same transaction.}
   \label{fig:soul_for_a_donut}
\end{figure}

From an economic argument alone, ethics (which is probably the only tool that can fight corruption - I argue that briefly a little bit later) is worth frequent examination.  Most firms are doing their best to define and implement anti-corruption policies, but many lack consistency, fail to collect feedback, and fail to recognize acts of integrity as well as acts of corruption.   If integrity is not actively recognized in some meaningful fashion, policies do little good.  As an example, the U.S. Congress has a Code of Ethics, but if members don't bother to follow the code, and Congress fails to self-police, then the code is just wasted paper (The readers can make "jokes" at the expense of Congress at their leisure!). 

There is considerable lack of enforcement worldwide often because of a subtle but important distinction between legal and ethical behavior.  One common theme often cited is that corruption is perceived differently in various parts of the world, and thus a universally accepted code of conduct would be difficult because cultural differences conspire against a uniform view of corruption.  This statement may be true in a legal context, but in an ethical context I argue the opposite; Ethics provides considerable clarity.  Virtue, duty, and utility are uniform principles regardless of culture - and that is likely why the emphasis on ethics is so important in engineering. \footnote{Most of this section is based on an article by Jeff Brown in the November 2004 ASCE News.}

Relevant to this discussion is the corporate mentality regarding morality.  In principle, a corporation is not a person, hence it cannot be a moral agent.  Because it is not a moral agent, corporations can neither be ethical nor unethical.  However, corporations are comprised of people and deal with people.
Because of these interactions with people, corporations are expected to behave morally, even though it is unenforceable (in the legal sense). 